% Source: https://artofproblemsolving.com/community/c3h1535233_area_of_an_octagon
% Stable identity: aops-topic-1535233-q01
Extend four alternate sides of the octagon to form a square; the octagon is the square with four corner triangles removed:

\begin{asy}
usepackage("amsmath");
size(5.5cm);
defaultpen(fontsize(10pt));
real s=6,t=s/sqrt(2),L=s+2*t;
draw((0,0)--(L,0)--(L,L)--(0,L)--cycle,dashed);
pair[] V={(t,0),(t+s,0),(L,t),(L,t+s),(t+s,L),(t,L),(0,t+s),(0,t)};
fill(V[0]--V[1]--V[2]--V[3]--V[4]--V[5]--V[6]--V[7]--cycle,paleblue);
draw(V[0]--V[1]--V[2]--V[3]--V[4]--V[5]--V[6]--V[7]--cycle,linewidth(1));
label("$6$",(L/2,0),N);
label("$6$",(0.55*t,0.55*t),NE);
label("$3\sqrt2$",(t/2,0),S);label("$3\sqrt2$",(0,t/2),W);
real q=0.7; draw((q,0)--(q,q)--(0,q));
label("\mbox{}",(min(currentpicture,user=true)+max(currentpicture,user=true))/2,black); // keeps the text after the figure black when LaTeX typesets the labels inline
\end{asy}

Each corner triangle is an isosceles right triangle with hypotenuse $6$ (a side of the octagon; the exterior angles of a regular octagon are $45^\circ$), so its legs are $\dfrac{6}{\sqrt2}=3\sqrt2$. The square has side $3\sqrt2+6+3\sqrt2=6+6\sqrt2$. Therefore
\[
\text{area}=(6+6\sqrt2)^2-4\cdot\dfrac{1}{2}(3\sqrt2)^2=(36+72\sqrt2+72)-36=\boxed{72+72\sqrt2}.
\]
(This agrees with the formula the poster mentions: the distance from the centre to a side is half the side of the square, $3+3\sqrt2$, and $\dfrac{1}{2}\cdot48\cdot(3+3\sqrt2)=72+72\sqrt2$.)
